%!TEX root = ../hutbthesis_main.tex
% 设置英文摘要
\keywordsen{Musculoskeletal Control; Motion Imitation; Reinforcement Learning; Quantitative Evaluation; Kinesis; Reproducible Experiments.Biological Rationality Algorithm Optimization}
\begin{abstracten}

To address the issues with low mimetic accuracy, unstable dynamics, and
weak biomechanical feasibility in current musculoskeletal locomotion
simulations, this paper explores biologically precise movement control
strategies within the Kinesis framework. The approach centers on Deep
Reinforcement Learning, which drives the core of a multi-type reward
function. Alongside this, four penalties—position deviation, velocity
mismatch, biomechanical limits, and failure to reach the goal—push the
system toward greater accuracy and real-world applicability.Three main
metrics guide the evaluation: Mean Per-Joint Position Error (MPJPE),
temporal coherence, and functional performance. Testing focuses on
replicative fidelity, dynamic reliability, and physical robustness.
Results show that blending diverse incentives with stricter biological
constraints not only reduces skeletal errors but also prevents kinematic
breakdowns and encourages smoother, more lifelike movement. The data
confirm the method's effectiveness and edge in muscle-skeleton
excitation. These standards matter directly for fields such as assistive
robotics, virtual avatar modeling, and mobile aid devices.

\end{abstracten}
